\documentclass{article}
\usepackage{ita2026}

\renewcommand{\ITAshorttitle}{Information Technology Applications (ITA)}
\renewcommand{\ITAshortauthor}{Juraj Suchán}

\makeatletter
%\def\input@path{{sections/}}
\makeatother

\begin{document}

\ITAjournalhead
  {International Journal\\Information Technology Applications (ITA)}
  {Volume XX, Number X, Feb. 2026}

\ITAtitle{Design and Implementation of a Workflow-Enabled Enterprise\\Document Management System with Legislative Compliance}
\ITAauthors{Juraj Suchán} 

\begin{ITAabstract}
This paper presents the design and implementation of an enterprise document and records management system with integrated workflow support. 
The proposed solution addresses the need for structured management of registry records, including classification, retention policies, disposal procedures, and auditability, 
while maintaining extensibility and modern architectural principles. 
The system is based on a microservice-oriented architecture implemented in .NET, utilizing a relational database for structured metadata storage and distributed components for messaging and caching. 

A dedicated workflow engine enables the definition and execution of configurable process models, including task, decision, and scripted nodes, allowing dynamic business logic execution and versioned workflow definitions. 
The data model reflects hierarchical classification structures and regulatory requirements related to records lifecycle management. 
The implementation emphasizes modularity, traceability, and compliance with records management principles \cite{iso15489,moreq2010}.

The resulting system demonstrates that enterprise-grade document and records management functionality can be combined with flexible workflow execution within a unified architecture, 
while preserving scalability and maintainability.
\end{ITAabstract}

\ITAkeywords{\textit{document management system, workflow engine, legislative compliance}}
\ITAacmccs{Software and its engineering $\rightarrow$ Software creation and management $\rightarrow$ Software development techniques}

%!TEX root = ../main.tex
\section*{Introduction}

Enterprise document and records management systems (DMS/RMS) are fundamental components of contemporary information infrastructures. 
They ensure structured storage, controlled access, traceability, and lifecycle management of organizational documents. 
In regulated environments, such systems must additionally support compliance with legal and normative requirements related to classification schemes, retention periods, disposal procedures, and auditability.

While traditional document management platforms primarily address storage and retrieval functionality, advanced process coordination is often delegated to separate workflow or business process management systems. 
This architectural separation increases integration complexity, introduces data redundancy, and may result in inconsistencies in document lifecycle control. 
Consequently, organizations face challenges in maintaining a coherent link between records management principles and operational process execution. 
Existing standards such as ISO 15489 and MoReq2010 define conceptual and functional requirements for records management~\cite{iso15489,moreq2010}, yet their practical integration with configurable workflow engines is frequently left to system-specific implementations.

This paper presents the design and implementation of a prototype enterprise document and records management system with integrated workflow support. 
The proposed solution adopts a unified architectural approach in which structured records modeling and configurable workflow execution are treated as mutually dependent components rather than isolated subsystems. 
The system is implemented using modern architectural principles, emphasizing modularity, scalability, and compliance-oriented data modeling.

The primary contribution of this work lies in: 
(i) the integration of a versioned and configurable workflow engine directly into a structured records management model, 
(ii) the alignment of workflow execution mechanisms with hierarchical classification structures and retention policies, and 
(iii) the architectural design of a modular system that supports auditability and extensibility without sacrificing maintainability.

The objective of this paper is to describe the architectural design, data model, workflow engine structure, and compliance mechanisms of the proposed solution, 
and to demonstrate how these elements form a coherent and extensible enterprise-ready system.
\setcounter{section}{0}

\input{sections/02-related-work}
\input{sections/03-architecture}
\input{sections/04-workflow-engine}
\input{sections/05-compliance}
\input{sections/06-discussion}
%!TEX root = ../main.tex
\section{Conclusion}

This paper presented the architectural design and implementation of a workflow-enabled enterprise document and records management system that integrates compliance-oriented modeling directly into its core domain structure. 
Rather than treating workflow orchestration and records management as independent subsystems, the proposed solution unifies classification hierarchies, lifecycle attributes, configurable metadata, and versioned workflow execution within a single coherent architecture.

The presented model demonstrates that regulatory alignment, structural consistency, and process flexibility can coexist without introducing architectural fragmentation. 
By separating workflow definitions from runtime execution state, enforcing schema-driven metadata validation, and embedding retention logic into classification entities, the system establishes a compliance-first foundation while preserving configurability and extensibility.

The architectural approach further illustrates how modern service-oriented design, asynchronous processing, and container-based deployment can support scalable and maintainable enterprise registry solutions. 
The hybrid metadata persistence strategy and versioned workflow model contribute to both adaptability and auditability, which are essential characteristics of regulated information systems.

Future work will focus on automated lifecycle enforcement mechanisms, extended compliance validation layers, and empirical performance evaluation under large-scale distributed workloads. 
Additional research may also explore generalization of the proposed model to other regulated domains requiring traceable document-centric workflows.


\clearpage

\bibliographystyle{plainnat}
\bibliography{refs}

\end{document}


